\chapter{Fazit und Ausblick}
\label{ch:conclusion}

Diese Arbeit hat \DO als einen primär soziokulturellen Transformationsprozess analysiert und dessen Erfolg von menschlichen und organisationalen Faktoren abhängig gemacht.
Das Fazit bündelt die zentralen Ergebnisse, beantwortet die Forschungsfrage und zeigt die Implikationen für Wissenschaft und Praxis auf.

\section{Hauptergebnisse}
Die Analyse der wissenschaftlichen Literatur hat zu folgenden zentralen Schlussfolgerungen geführt:

\begin{enumerate}
	\item \DO ist primär eine kulturelle und nur sekundär eine technologische Transformation. Der Erfolg hängt weniger von der Werkzeugauswahl ab als von der Etablierung einer kollaborativen Denk"= und Arbeitsweise.
	\item Eine kollaborative Vertrauenskultur, unterstützt durch transformative Führung, ist der entscheidende Erfolgsfaktor. Vertrauen, psychologische Sicherheit und offene Kommunikation bilden das Fundament einer funktionierenden \DO-Praxis.
	\item Traditionelle organisationale Silos und die Angst vor dem Unbekannten stellen die größten Widerstände dar. Starre Hierarchien und funktionale Abteilungsgrenzen unterbinden die für \DO notwendige Agilität und Zusammenarbeit.
\end{enumerate}

Basierend auf diesen Erkenntnissen kann die eingangs gestellte Forschungsfrage \enquote{Welche sozialen und kulturellen Faktoren fördern oder hemmen die \DO-Transformation in Software-Unternehmen?} wie folgt beantwortet werden:

Die entscheidenden sozialen und kulturellen Erfolgsfaktoren sind Vertrauen, offene Kommunikation, geteilte Verantwortung und eine transformative Führung, die für ein gemeinsames Verständnis von \DO sorgt.
Die größten Widerstände entstehen aus traditionellen Hierarchien, Silodenken und einem kulturellen Spannungsfeld zwischen Stabilitäts"= und Innovationszielen.
Der Aufbau einer \DO-Kultur erfordert daher gezielte Maßnahmen im Change"=Management, die diese sozialen Dynamiken adressieren und eine nachhaltige kulturelle Transformation ermöglichen.
Arbeitssoziologisch betrachtet stellt \DO ein grundlegend neues Konzept dar, das die Art und Weise der Zusammenarbeit in der Informationstechnik reformieren kann.

\section{Bedeutung für Forschung und Praxis}
Die Ergebnisse haben eine hohe Relevanz für die unternehmerische Praxis und die wissenschaftliche Forschung.

Für die Praxis und die Organisationsentwicklung bieten die Befunde eine klare Handlungsorientierung.
Sie bestätigen, dass der Fokus von Transformationsbemühungen auf Kultur, Führung und Kommunikation liegen muss.
Eine reine Technologie"=Implementierung ohne begleitende kulturelle Maßnahmen ist daher unzureichend.

\DO ermöglicht in der soziologischen Forschung die empirische Untersuchung einer neuen Arbeitsform, die die Dynamik in interdisziplinären Teams und den Wandel von Machtstrukturen in technologiegetriebenen Organisationen widerspiegelt.
Es dient damit als paradigmatischer Fall für die Transformation der Arbeit im 21. Jahrhundert und zeigt, wie menschliche Systeme zusammen mit technischen Systemen neu ausgerichtet werden müssen, um strategische Ziele zu erreichen.

\section{Offene Fragen und Ausblick}
Zukünftige Untersuchungen könnten sich auf spezifischere Aspekte konzentrieren:
\begin{itemize}
	\item Die Integration von Sicherheitsaspekten in die \DO-Kultur (Dev\-Sec\-Ops) und die damit verbundenen sozialen und prozessualen Herausforderungen.
	\item Der Einfluss von Künstlicher Intelligenz (KI) auf \DO-Prozesse (AI\-Ops) und die daraus resultierenden Veränderungen für die Rollen und Kompetenzen der Mitarbeiter.
	\item Längsschnittstudien, die den kulturellen Wandel in Organisationen über mehrere Jahre begleiten, um die langfristigen Effekte von \DO-Trans"-for"-ma"-tio"-nen zu untersuchen.
\end{itemize}

Die kontinuierliche Auseinandersetzung mit diesen Themen ist entscheidend, da sich \DO weiterentwickelt und ein zentrales Paradigma für die Zukunftsfähigkeit von Unternehmen im digitalen Zeitalter bleibt.